
\chapter{Kalman Filter}

% ==========================================================
\section{The filtering problem}

% --- Add this at the end of Section 1 (after "The filtering problem in the course notation") ---

\paragraph{Motivation and intuition: what filtering is trying to achieve}
Think of a pair \((\theta_t,\xi_t)\), which we will name \emph{hidden state} and \emph{measurement}, respectively.
At each time \(t\), the state \(\theta_t\) contains what we would really like to know (e.g.\ position/velocity, a temperature field coefficient, a parameter vector), but we do not observe it directly.
Instead we only observe \(\xi_0,\dots,\xi_t\), which are corrupted by noise and often contain only partial information about \(\theta_t\).
The \emph{filtering problem} is to turn the data stream \(\xi_0,\xi_1,\dots\) into a running estimate of the current state.

Mathematically, ``using all information available up to time \(t\)'' is expressed by restricting ourselves to \(\F_t^\xi\)-measurable estimators.
Among all such estimators, the conditional expectation
\[
m(t)=\E[\theta_t\mid \F_t^\xi]
\]
is optimal in the mean-square sense: it is the unique (a.s.) estimator that minimizes the expected squared error
\(\E[\|\theta_t-\hat\theta\|^2]\) over all \(\F_t^\xi\)-measurable \(\hat\theta\).
So the filter does not guess \(\theta_t\); it computes the best possible estimate given exactly what has been observed so far.

Intuitively, filtering is a repeated \emph{predict--correct} cycle.
From the model dynamics we can \emph{predict} what \(\theta_{t+1}\) should look like given \(\F_t^\xi\) (a ``prior'' belief).
Then the new measurement \(\xi_{t+1}\) arrives and we \emph{correct} this prediction by comparing \(\xi_{t+1}\) to what the model expected to see.
This comparison is encoded by the \emph{innovation} (or residual)
\[
\nu_{t+1}=\xi_{t+1}-\E[\xi_{t+1}\mid \F_t^\xi],
\]
which measures the genuinely new information in \(\xi_{t+1}\) that was not predictable from past data.
The conditional covariance
\[
Y(t)=\Cov(\theta_t\mid \F_t^\xi)
\]
quantifies the remaining uncertainty: large \(Y(t)\) means the observations so far do not pin down the state well, while small \(Y(t)\) means the filter is confident.
In the (conditionally) Gaussian setting of these notes, keeping track of the pair \((m(t),Y(t))\) is enough to describe the entire conditional law of \(\theta_t\mid \F_t^\xi\), so the infinite-dimensional problem ``update a distribution'' reduces to a finite-dimensional recursion for mean and covariance. The key idea is that if, conditional on the observation
history, the next-step pair \((\theta_{t+1},\xi_{t+1})\) is Gaussian, then the full filter reduces to
updating the conditional mean \(m(t)=\E[\theta_t\mid \F_t^\xi]\) and conditional covariance
\(Y(t)=\Cov(\theta_t\mid \F_t^\xi)\) via the normal-correlation (Gaussian conditioning) formulas.
A pseudo-inverse is used so the update still makes sense when conditional observation covariances are singular.


\subsection{State, observations, and information}
Let \((\Omega,\F,\mathbb{P})\) be a probability space.
We consider two processes:
\begin{itemize}
  \item an \emph{unobserved state} \(\theta_t\in\R^k\),
  \item an \emph{observation/measurement} \(\xi_t\in\R^\ell\),
\end{itemize}
indexed by discrete time \(t=0,1,2,\dots\).
The information available up to time \(t\) is the observation filtration
\[
\F_t^\xi \coloneqq \sigma(\xi_0,\xi_1,\dots,\xi_t).
\]
The goal of filtering is to estimate \(\theta_t\) from \(\F_t^\xi\).

\subsection{The model}
% --- Add this paragraph right after Assumption 2.1 (Course model for (\theta_t,\xi_t)) ---

Consider a linear system driven by Gaussian noises, but allow the \emph{coefficients}
to be random and adapted to the observation history.

For \(t\ge 0\),
\begin{align}
\theta_{t+1} &= a_0(t) + a_1(t)\theta_t + b_1(t)\,\varepsilon_1(t+1)+b_2(t)\,\varepsilon_2(t+1), \label{eq:theta}\\
\xi_{t+1} &= A_0(t) + A_1(t)\theta_t + B_1(t)\,\varepsilon_1(t+1)+B_2(t)\,\varepsilon_2(t+1). \label{eq:xi}
\end{align}
Here:
\begin{itemize}
  \item \(\varepsilon_1(t)\in\R^{r_1}\) and \(\varepsilon_2(t)\in\R^{r_2}\) are i.i.d.\ standard Gaussian vectors,
        independent across time and mutually independent;
  \item for each \(t\), the coefficients \(a_0(t),a_1(t),b_1(t),b_2(t),A_0(t),A_1(t),B_1(t),B_2(t)\)
        are \(\F_t^\xi\)-measurable (i.e.\ can depend on past observations).
\end{itemize}

\paragraph{Why the model takes the form \eqref{eq:theta}--\eqref{eq:xi} (intuition).}
The two equations \eqref{eq:theta}--\eqref{eq:xi} are a concise way to encode the standard ingredients of a
\emph{state--space model}: (i) a \emph{dynamical law} describing how the hidden state evolves from time \(t\) to \(t+1\),
and (ii) an \emph{observation law} describing how the measurement is generated from the current state.
The evolution equation \eqref{eq:theta} says: starting from \(\theta_t\), we apply a linear ``propagation''
\(a_1(t)\theta_t\) plus a drift/offset \(a_0(t)\), and then we add random perturbations
\(b_1(t)\varepsilon_1(t+1)+b_2(t)\varepsilon_2(t+1)\) which model unmodeled dynamics (process noise).
The measurement equation \eqref{eq:xi} says: the sensor sees a linear function \(A_1(t)\theta_t\) of the state
(with possible bias \(A_0(t)\)), and it is corrupted by noise
\(B_1(t)\varepsilon_1(t+1)+B_2(t)\varepsilon_2(t+1)\).
Allowing the coefficients to be \(\F_t^\xi\)-measurable means the system can be \emph{adaptive}:
at time \(t\) the model parameters may depend on past observations (e.g.\ regime switching, time-varying calibration,
feedback control), but crucially the dependence is only on information already available at time \(t\).
Finally, the use of Gaussian noises and linear combinations is exactly what preserves (conditional) Gaussianity,
which is why the filter can be propagated using only conditional means and covariances.

\begin{remark}[Why this is ``more general'' than the standard Kalman filter]
In the classical linear-Gaussian state space model the matrices are deterministic (or at least independent of the observations).
Here, they may depend on the observation history (e.g.\ switching regimes, adaptive control, random coefficients).
Nevertheless, \eqref{eq:theta}--\eqref{eq:xi} remain \emph{linear} in the hidden state and the driving Gaussian noises,
so Gaussian conditioning still applies.
\end{remark}

\begin{remark}[Optimal (mean-square) estimate]
A basic fact from conditional expectation is:
\[
m(t)\coloneqq \E[\theta_t\mid \F_t^\xi]
\quad\text{minimizes}\quad
\E\bigl[\|\theta_t-\hat\theta\|^2\bigr]
\ \text{over all }\ \F_t^\xi\text{-measurable }\hat\theta.
\]
The associated conditional error covariance is
\[
Y(t)\coloneqq \Cov(\theta_t\mid \F_t^\xi)
= \E\!\left[(\theta_t-m(t))(\theta_t-m(t))^\transpose\,\middle|\,\F_t^\xi\right].
\]
In the classical Kalman filter, \((m(t),Y(t))\) evolve by a simple recursion.
\end{remark}


% ==========================================================
\section{Conditionally Gaussian sequences (Lipster--Shiryayev framework)}

\paragraph{Intuition and motivation (why conditional Gaussianity matters).}
Filtering is, at its core, a problem of updating a \emph{conditional distribution}:
after observing \(\xi_0,\dots,\xi_t\), our knowledge about the hidden state \(\theta_t\) is encoded by the law of
\(\theta_t\mid \F_t^\xi\).
In a general nonlinear/non-Gaussian model this conditional law is an infinite-dimensional object (a full density or measure),
so an ``exact'' filter would require propagating an entire function over time (which is why particle filters exist).
The point of Theorem~\ref{thm:cond_gaussian} is that for the model \eqref{eq:theta}--\eqref{eq:xi} this complexity collapses:
\emph{given the observation history, the state remains Gaussian}.
A Gaussian law is completely determined by its mean and covariance, so instead of propagating an entire distribution we only
need to propagate the finite-dimensional pair \(\bigl(m(t),Y(t)\bigr)\).
This is precisely the structural reason the Kalman filter is computationally simple.

Why is conditional Gaussianity plausible here?
Start from a Gaussian prior (or, more generally, assume \(\theta_0\) has a Gaussian conditional law given \(\F_0^\xi\)).
Then each step takes linear combinations of the previous state \(\theta_t\) and fresh Gaussian noises
\(\varepsilon_1(t+1),\varepsilon_2(t+1)\).
Conditional on \(\F_t^\xi\), the coefficients \(a_0(t),a_1(t),\dots\) are just \emph{known numbers/matrices}
(because they are \(\F_t^\xi\)-measurable), so \(\theta_{t+1}\) becomes an affine transformation of Gaussian variables,
hence Gaussian.
Moreover, the pair \((\theta_{t+1},\xi_{t+1})\) is jointly Gaussian given \(\F_t^\xi\), because \(\xi_{t+1}\) is built from
the same Gaussian ingredients in \eqref{eq:xi}.
Once we have joint (conditional) Gaussianity, the update step
\(\theta_{t+1}\mid \F_{t+1}^\xi=\sigma(\F_t^\xi,\xi_{t+1})\)
is obtained by the standard Gaussian conditioning (normal-correlation) formulas,
which is exactly what produces the recursion for \(m(t+1)\) and \(Y(t+1)\).
So Theorem~\ref{thm:cond_gaussian} is the bridge that justifies replacing the abstract filtering problem
``compute \(\mathcal{L}(\theta_t\mid \F_t^\xi)\)'' by the concrete matrix recursion of the generalized Kalman filter.


Under mild square-integrability and boundedness assumptions on the coefficients (as stated in your notes),
one can show that for each \(t\) the conditional distribution of \(\theta_t\) given \(\F_t^\xi\) is Gaussian a.s.
This is one of the main structural results used to justify that it is enough to propagate \((m(t),Y(t))\).

\begin{theorem}[Conditional Gaussianity: informal statement]
\label{thm:cond_gaussian}
Assume the coefficients in \eqref{eq:theta}--\eqref{eq:xi} are square-integrable and uniformly bounded in \(t\) in probability,
and \(\E(\|\theta_0\|^2+\|\xi_0\|^2)<\infty\).
Then for each \(t\ge 0\),
\[
\theta_t\mid \F_t^\xi \ \sim\ \mathcal{N}\bigl(m(t),\,Y(t)\bigr)
\quad\text{(almost surely)}.
\]
\end{theorem}

\begin{proof}[Proof (split into steps, following Lipster--Shiryayev)]
We prove by induction on $t$ that the conditional law of $\theta_t$ given $\F_t^\xi$ is (a.s.) Gaussian.

\medskip
\noindent\textbf{Step 0: Base case.}
We assume the initial condition is \emph{conditionally Gaussian} in the sense that
\[
\theta_0 \mid \F_0^\xi \sim \mathcal{N}(m(0),Y(0)) \quad \text{a.s.}
\]
(For instance, this holds if $(\theta_0,\xi_0)$ is jointly Gaussian; then $\theta_0\mid \xi_0$
is Gaussian by the normal-correlation formula.)

\medskip
\noindent\textbf{Step 1: Induction hypothesis.}
Fix $t\ge 0$ and assume
\[
\theta_t \mid \F_t^\xi \sim \mathcal{N}(m(t),Y(t)) \quad \text{a.s.}
\]
We show that then $\theta_{t+1}\mid \F_{t+1}^\xi$ is also Gaussian.

\medskip
\noindent\textbf{Step 2: Write $(\theta_{t+1},\xi_{t+1})$ as an affine function of Gaussian variables.}
Recall the model
\begin{align*}
\theta_{t+1} &= a_0(t) + a_1(t)\theta_t + b_1(t)\varepsilon_1(t+1)+b_2(t)\varepsilon_2(t+1),\\
\xi_{t+1} &= A_0(t) + A_1(t)\theta_t + B_1(t)\varepsilon_1(t+1)+B_2(t)\varepsilon_2(t+1),
\end{align*}
where all coefficients are $\F_t^\xi$-measurable and $\varepsilon_1,\varepsilon_2$ are i.i.d.\ standard Gaussian,
independent of $\F_t^\xi$.

Define the stacked Gaussian vector
\[
U_{t+1}\coloneqq
\begin{pmatrix}
\theta_t\\ \varepsilon_1(t+1)\\ \varepsilon_2(t+1)
\end{pmatrix}.
\]
Conditional on $\F_t^\xi$, the vector $U_{t+1}$ is Gaussian because:
(i) $\theta_t\mid \F_t^\xi$ is Gaussian by the induction hypothesis, and
(ii) $\varepsilon_1(t+1),\varepsilon_2(t+1)$ are Gaussian and independent of $\F_t^\xi$.
Moreover,
\[
\E[U_{t+1}\mid \F_t^\xi]=
\begin{pmatrix}
m(t)\\ 0\\ 0
\end{pmatrix},
\qquad
\Cov(U_{t+1}\mid \F_t^\xi)=
\begin{pmatrix}
Y(t) & 0 & 0\\
0 & I & 0\\
0 & 0 & I
\end{pmatrix}.
\]

Now introduce the $\F_t^\xi$-measurable affine map
\[
\begin{pmatrix}
\theta_{t+1}\\ \xi_{t+1}
\end{pmatrix}
=
\underbrace{\begin{pmatrix} a_0(t)\\ A_0(t)\end{pmatrix}}_{=:c(t)}
+
\underbrace{\begin{pmatrix}
a_1(t) & b_1(t) & b_2(t)\\
A_1(t) & B_1(t) & B_2(t)
\end{pmatrix}}_{=:M(t)}
\begin{pmatrix}
\theta_t\\ \varepsilon_1(t+1)\\ \varepsilon_2(t+1)
\end{pmatrix}
=
c(t)+M(t)U_{t+1}.
\]
Since $U_{t+1}\mid \F_t^\xi$ is Gaussian and $(c(t),M(t))$ are $\F_t^\xi$-measurable (hence ``constants'' under the conditioning),
it follows that the pair $(\theta_{t+1},\xi_{t+1})\mid \F_t^\xi$ is \emph{jointly Gaussian}.

\medskip
\noindent\textbf{Step 3: Identify the conditional mean/covariance blocks.}
From the affine representation above (or by direct expansion using independence of the noises), we obtain
\[
\E[\theta_{t+1}\mid \F_t^\xi]=a_0(t)+a_1(t)m(t),\qquad
\E[\xi_{t+1}\mid \F_t^\xi]=A_0(t)+A_1(t)m(t),
\]
and the conditional covariance blocks
\begin{align*}
d_{11}^\xi(t)&\coloneqq \Cov(\theta_{t+1}\mid \F_t^\xi)
= a_1(t)Y(t)a_1(t)^\transpose + b_1(t)b_1(t)^\transpose + b_2(t)b_2(t)^\transpose,\\
d_{22}^\xi(t)&\coloneqq \Cov(\xi_{t+1}\mid \F_t^\xi)
= A_1(t)Y(t)A_1(t)^\transpose + B_1(t)B_1(t)^\transpose + B_2(t)B_2(t)^\transpose,\\
d_{12}^\xi(t)&\coloneqq \Cov(\theta_{t+1},\xi_{t+1}\mid \F_t^\xi)
= a_1(t)Y(t)A_1(t)^\transpose + b_1(t)B_1(t)^\transpose + b_2(t)B_2(t)^\transpose.
\end{align*}

\medskip
\noindent\textbf{Step 4: Condition on the new observation (normal correlation).}
Because $(\theta_{t+1},\xi_{t+1})\mid \F_t^\xi$ is jointly Gaussian, we can condition on $\xi_{t+1}$ using the
Gaussian regression (normal-correlation) formula. Since
\(
\F_{t+1}^\xi=\sigma(\F_t^\xi,\xi_{t+1}),
\)
this is exactly conditioning on the available information at time $t+1$.

Let $\nu_{t+1}\coloneqq \xi_{t+1}-\E[\xi_{t+1}\mid \F_t^\xi]$ be the innovation.
Then (allowing $d_{22}^\xi(t)$ to be singular) the conditional law of $\theta_{t+1}$ given $\F_{t+1}^\xi$ is Gaussian with
\begin{align*}
m(t+1)
&\coloneqq \E[\theta_{t+1}\mid \F_{t+1}^\xi]
= \E[\theta_{t+1}\mid \F_t^\xi] + d_{12}^\xi(t)\bigl(d_{22}^\xi(t)\bigr)^{+}\,\nu_{t+1},\\
Y(t+1)
&\coloneqq \Cov(\theta_{t+1}\mid \F_{t+1}^\xi)
= d_{11}^\xi(t) - d_{12}^\xi(t)\bigl(d_{22}^\xi(t)\bigr)^{+} d_{12}^\xi(t)^\transpose.
\end{align*}
In particular, $\theta_{t+1}\mid \F_{t+1}^\xi$ is (a.s.) Gaussian.

\medskip
\noindent\textbf{Step 5: Conclude the induction.}
Steps 2--4 show that the Gaussianity property propagates from time $t$ to time $t+1$.
Together with the base case (Step 0), this proves that for every $t\ge 0$ the conditional distribution
$\theta_t\mid \F_t^\xi$ is (a.s.) Gaussian.
\end{proof}


\begin{remark}
The paper by Lipster--Shiryayev develops this point of view systematically and calls such processes
\emph{conditionally Gaussian sequences}.
The entire filtering problem then reduces to finding recursion relations for \(m(t)\) and \(Y(t)\).
\end{remark}

% ==========================================================
\section{Gaussian conditioning and the pseudo-inverse}

We learned that under certain conditions we know that $(\theta_{t+1},\xi{t+1})$ are jointly Gaussian, the only thing remaining is to compute them. The next theorem shows how we can perform this step.
The only twist is that the conditional covariance of \(\xi\) may be singular; then one uses a pseudo-inverse.

\begin{theorem}[Normal correlation / Gaussian regression with pseudo-inverse]
\label{thm:gaussian_regression}
Let \(X\in\R^k\), \(Z\in\R^\ell\) be jointly Gaussian with mean
\(\E[X]=m_X\), \(\E[Z]=m_Z\) and block covariance
\[
\Sigma=
\begin{pmatrix}
\Sigma_{XX} & \Sigma_{XZ}\\
\Sigma_{ZX} & \Sigma_{ZZ}
\end{pmatrix},
\qquad \Sigma_{ZZ}\ \text{symmetric, nonnegative definite (possibly singular)}.
\]
Then
\begin{align*}
\E[X\mid Z] &= m_X + \Sigma_{XZ}\,\Sigma_{ZZ}\pinv\,(Z-m_Z),\\
\Cov(X\mid Z) &= \Sigma_{XX} - \Sigma_{XZ}\,\Sigma_{ZZ}\pinv\,\Sigma_{ZX}.
\end{align*}
If \(\Sigma_{ZZ}\) is invertible, then \(\Sigma_{ZZ}\pinv=\Sigma_{ZZ}^{-1}\) and we recover the familiar formulas.
\end{theorem}

\begin{remark}[Interpretation]
\(\Sigma_{XZ}\Sigma_{ZZ}\pinv\) is the optimal linear gain mapping a measurement residual \((Z-m_Z)\)
to a correction of the mean of \(X\).
Using a pseudo-inverse allows the formula to remain valid if some linear combinations of the measurement are noiseless
or redundant (hence \(\Sigma_{ZZ}\) becomes singular).
\end{remark}

\paragraph{Motivation and intuition.}
Theorem~\ref{thm:gaussian_regression} is the mathematical engine behind the Kalman update.
In filtering, after time \(t\) we have a \emph{prior} joint (conditional) Gaussian law for the pair
\((X,Z)=(\theta_{t+1},\xi_{t+1})\) given the past \(\F_t^\xi\); when the new measurement \(Z\) is revealed,
we need the \emph{posterior} law of \(X\mid Z\).
For Gaussian vectors, this posterior stays Gaussian, and its mean is obtained by the best linear
least-squares correction of the prior mean: one takes the residual \((Z-m_Z)\) and maps it to a correction of \(X\)
through a gain matrix \( \Sigma_{XZ}\Sigma_{ZZ}\pinv \).
Geometrically, this is the orthogonal projection of \(X\) onto the affine space generated by \(Z\)
(in \(L^2\)), so only the part of \(X\) that is correlated with \(Z\) is ``explained'' by the observation;
the covariance update subtracts precisely that explained part.
The pseudo-inverse appears because \(\Sigma_{ZZ}\) may fail to be invertible (e.g.\ redundant sensors,
noiseless components, or linear dependencies), but the regression/projection viewpoint still makes sense:
\(\Sigma_{ZZ}\pinv\) selects the minimum-norm gain consistent with the information contained in \(Z\),
and the formulas remain valid for any positive semidefinite \(\Sigma_{ZZ}\).

\paragraph{Relation to Theorem~\ref{thm:cond_gaussian} (Conditional Gaussianity).}
Theorem~\ref{thm:cond_gaussian} is the \emph{structural} statement that makes Theorem~\ref{thm:gaussian_regression} applicable at every time step:
it tells us that, under the model \eqref{eq:theta}--\eqref{eq:xi}, the relevant random objects remain Gaussian after conditioning on the
available information. More precisely, Theorem~\ref{thm:cond_gaussian} implies that given \(\F_t^\xi\) the pair
\((\theta_{t+1},\xi_{t+1})\) is \emph{jointly Gaussian} (because it is an affine function of the conditionally Gaussian \(\theta_t\)
and fresh Gaussian noises with \(\F_t^\xi\)-measurable coefficients). Once this joint (conditional) Gaussianity holds,
updating from time \(t\) to \(t+1\) is exactly the problem of computing
\(\theta_{t+1}\mid \sigma(\F_t^\xi,\xi_{t+1})\), i.e.\ conditioning one component of a jointly Gaussian vector on the other.
That is precisely what Theorem~\ref{thm:gaussian_regression} provides, producing the Kalman-type ``mean + gain $\times$ innovation'' update and the corresponding covariance reduction.


\begin{lemma}[Conditional Gaussian regression]\label{lem:cond_gauss_regression}
Let $(\Omega,\mathcal F,\mathbb P)$ be a probability space and let $\mathcal G\subset\mathcal F$ be a sub-$\sigma$-algebra.
Let $X:\Omega\to\mathbb R^{n}$ and $Z:\Omega\to\mathbb R^{m}$ be square-integrable random vectors.
Assume that, $\mathbb P$-a.s., the conditional law of $(X,Z)$ given $\mathcal G$ is (multivariate) Gaussian, i.e.,
there exist $\mathcal G$-measurable random vectors $\mu_X\in\mathbb R^{n}$, $\mu_Z\in\mathbb R^{m}$ and
$\mathcal G$-measurable random matrices
\[
\Sigma_{XX}\in\mathbb R^{n\times n},\qquad
\Sigma_{XZ}\in\mathbb R^{n\times m},\qquad
\Sigma_{ZZ}\in\mathbb R^{m\times m},
\]
such that for $\mathbb P$-a.e.\ $\omega$,
\[
\mathcal L\big((X,Z)\mid \mathcal G\big)(\omega)
=\mathcal N\!\left(
\begin{pmatrix}\mu_X(\omega)\\ \mu_Z(\omega)\end{pmatrix},
\begin{pmatrix}\Sigma_{XX}(\omega) & \Sigma_{XZ}(\omega)\\
\Sigma_{XZ}(\omega)^\top & \Sigma_{ZZ}(\omega)
\end{pmatrix}
\right).
\]
Then, with $\mu_X=\mathbb E[X\mid \mathcal G]$, $\mu_Z=\mathbb E[Z\mid \mathcal G]$,
$\Sigma_{XZ}=\mathrm{Cov}(X,Z\mid \mathcal G)$ and $\Sigma_{ZZ}=\mathrm{Cov}(Z\mid \mathcal G)$, we have
\begin{equation}\label{eq:cond_gauss_regression}
\mathbb E\big[X\mid \mathcal G, Z\big]
=
\mu_X + \Sigma_{XZ}\,\Sigma_{ZZ}^{+}\,(Z-\mu_Z)
\qquad \mathbb P\text{-a.s.},
\end{equation}
where $\Sigma_{ZZ}^{+}$ denotes the Moore--Penrose pseudoinverse.

If $\Sigma_{ZZ}$ is $\mathbb P$-a.s.\ invertible, then $\Sigma_{ZZ}^{+}=\Sigma_{ZZ}^{-1}$ and
\eqref{eq:cond_gauss_regression} becomes the familiar gain-times-innovation formula.
\end{lemma}

\begin{proof}
Fix $t$-like notation $\mathcal G$ and define $Y:=(X,Z)$.
Let $\pi_X:\mathbb R^{n+m}\to\mathbb R^{n}$ and $\pi_Z:\mathbb R^{n+m}\to\mathbb R^{m}$
be the coordinate projections, so that $X=\pi_X(Y)$ and $Z=\pi_Z(Y)$.

\medskip
\noindent\textbf{Step 1 (Work under the conditional probability).}
Let $\mathbb P(\cdot\mid \mathcal G)(\omega)$ be a regular conditional probability kernel for $\mathcal G$.
For $\mathbb P$-a.e.\ $\omega$, set
\[
\mathbb P^\omega(A):=\mathbb P(A\mid \mathcal G)(\omega),\qquad A\in\mathcal F.
\]
By assumption, under $\mathbb P^\omega$ the vector $Y=(X,Z)$ is an \emph{ordinary} Gaussian vector in
$\mathbb R^{n+m}$ with mean and covariance given by the $\mathcal G$-measurable blocks
$\mu_X(\omega),\mu_Z(\omega)$ and $\Sigma_{ij}(\omega)$.

\medskip
\noindent\textbf{Step 2 (Apply unconditional Gaussian regression under $\mathbb P^\omega$).}
For each such $\omega$, apply the standard Gaussian regression formula (Theorem~4.5 )
to the jointly Gaussian pair $(X,Z)$ under the probability measure $\mathbb P^\omega$:
\begin{equation}\label{eq:regression_under_Pomega}
\mathbb E_{\mathbb P^\omega}[X\mid Z]
=
\mathbb E_{\mathbb P^\omega}[X]
+
\mathrm{Cov}_{\mathbb P^\omega}(X,Z)\;
\mathrm{Cov}_{\mathbb P^\omega}(Z)^{+}\,
\bigl(Z-\mathbb E_{\mathbb P^\omega}[Z]\bigr).
\end{equation}
By the definition of the conditional mean/covariance blocks, we have for $\mathbb P$-a.e.\ $\omega$:
\[
\mathbb E_{\mathbb P^\omega}[X]=\mu_X(\omega),\quad
\mathbb E_{\mathbb P^\omega}[Z]=\mu_Z(\omega),\quad
\mathrm{Cov}_{\mathbb P^\omega}(X,Z)=\Sigma_{XZ}(\omega),\quad
\mathrm{Cov}_{\mathbb P^\omega}(Z)=\Sigma_{ZZ}(\omega).
\]
Substituting these identities into \eqref{eq:regression_under_Pomega} yields, for $\mathbb P$-a.e.\ $\omega$,
\begin{equation}\label{eq:pathwise_formula}
\mathbb E_{\mathbb P^\omega}[X\mid Z]
=
\mu_X(\omega)+\Sigma_{XZ}(\omega)\Sigma_{ZZ}(\omega)^{+}\bigl(Z-\mu_Z(\omega)\bigr).
\end{equation}

\medskip
\noindent\textbf{Step 3 (Identify $\mathbb E[X\mid \mathcal G,Z]$).}
By the defining property of regular conditional probabilities,
for any bounded Borel function $\varphi:\mathbb R^{m}\to\mathbb R$,
\[
\mathbb E\!\left[X\,\varphi(Z)\mid \mathcal G\right](\omega)
=
\int X(\omega')\,\varphi(Z(\omega'))\,\mathbb P^\omega(d\omega').
\]
Moreover, the conditional expectation $\mathbb E[X\mid \mathcal G,Z]$ is characterized as the
$\sigma(\mathcal G,Z)$-measurable random vector $U$ such that for all bounded
$\sigma(\mathcal G,Z)$-measurable scalars $H$,
\[
\mathbb E[XH]=\mathbb E[UH].
\]
Using the kernel $\mathbb P^\omega$ and taking $H=h(\omega)\,\varphi(Z)$ with bounded
$\mathcal G$-measurable $h$ and bounded Borel $\varphi$, one checks that
\[
\mathbb E[X\mid \mathcal G,Z](\omega)=\mathbb E_{\mathbb P^\omega}[X\mid Z]
\qquad\text{for }\mathbb P\text{-a.e.\ }\omega.
\]
Combining this identification with the pathwise formula \eqref{eq:pathwise_formula} gives
\eqref{eq:cond_gauss_regression}.

\medskip
Finally, when $\Sigma_{ZZ}$ is invertible, the pseudoinverse equals the inverse and the stated simplification follows.
\end{proof}





% ==========================================================
\section{Deriving the recursive filter}

Fix \(t\ge 0\) and suppose \(m(t)\) and \(Y(t)\) are known.
We derive \(m(t+1)\) and \(Y(t+1)\) from the model \eqref{eq:theta}--\eqref{eq:xi}.

\subsection{Prediction (conditional means)}
Taking \(\E[\cdot\mid \F_t^\xi]\) and using that \(\varepsilon_1(t+1),\varepsilon_2(t+1)\) are independent of \(\F_t^\xi\)
with zero mean, we get the \emph{one-step predictions}
\begin{align}
\E[\theta_{t+1}\mid \F_t^\xi] &= a_0(t) + a_1(t)m(t), \label{eq:predtheta}\\
\E[\xi_{t+1}\mid \F_t^\xi] &= A_0(t) + A_1(t)m(t). \label{eq:predxi}
\end{align}
Define the \emph{innovation} (measurement residual)
\[
\nu_{t+1}\coloneqq \xi_{t+1}-\E[\xi_{t+1}\mid \F_t^\xi]
= \xi_{t+1}-A_0(t)-A_1(t)m(t).
\]

\subsection{Prediction (conditional covariances)}
Write the centered ``prediction errors''
\begin{align*}
\theta_{t+1}-\E[\theta_{t+1}\mid \F_t^\xi]
&= a_1(t)\bigl(\theta_t-m(t)\bigr)
  + b_1(t)\varepsilon_1(t+1)+b_2(t)\varepsilon_2(t+1),\\
\xi_{t+1}-\E[\xi_{t+1}\mid \F_t^\xi]
&= A_1(t)\bigl(\theta_t-m(t)\bigr)
  + B_1(t)\varepsilon_1(t+1)+B_2(t)\varepsilon_2(t+1).
\end{align*}
Using independence and \(\Var(\varepsilon_i)=I\), the conditional covariance blocks of
\((\theta_{t+1},\xi_{t+1})\) given \(\F_t^\xi\) are
\begin{align}
d_{11}^\xi(t) &\coloneqq \Cov(\theta_{t+1}\mid \F_t^\xi)
= a_1(t)Y(t)a_1(t)^\transpose + b_1(t)b_1(t)^\transpose + b_2(t)b_2(t)^\transpose,\label{eq:d11}\\
d_{22}^\xi(t) &\coloneqq \Cov(\xi_{t+1}\mid \F_t^\xi)
= A_1(t)Y(t)A_1(t)^\transpose + B_1(t)B_1(t)^\transpose + B_2(t)B_2(t)^\transpose,\label{eq:d22}\\
d_{12}^\xi(t) &\coloneqq \Cov(\theta_{t+1},\xi_{t+1}\mid \F_t^\xi)
= a_1(t)Y(t)A_1(t)^\transpose + b_1(t)B_1(t)^\transpose + b_2(t)B_2(t)^\transpose.\label{eq:d12}
\end{align}

\subsection{Update (conditioning on the new measurement)}
Now \(\F_{t+1}^\xi = \sigma(\F_t^\xi,\xi_{t+1})\), so we update by conditioning the (conditionally) Gaussian vector
\((\theta_{t+1},\xi_{t+1})\) on \(\xi_{t+1}\).
Applying the Gaussian regression theorem with pseudo-inverse yields:

\begin{proposition}[Generalized Kalman update]
Define the (possibly random) \emph{Kalman gain}
\[
K(t)\coloneqq d_{12}^\xi(t)\, \bigl(d_{22}^\xi(t)\bigr)\pinv.
\]
Then the filtered mean and covariance satisfy
\begin{align}
m(t+1)
&= \E[\theta_{t+1}\mid \F_{t+1}^\xi] \nonumber\\
&= a_0(t)+a_1(t)m(t)\;+\;K(t)\,\nu_{t+1}, \label{eq:update_m}\\[2mm]
Y(t+1)
&= \Cov(\theta_{t+1}\mid \F_{t+1}^\xi)
= d_{11}^\xi(t) - K(t)\,d_{22}^\xi(t)\,K(t)^\transpose. \label{eq:update_Y}
\end{align}
Equivalently, using the blocks directly,
\[
Y(t+1)= d_{11}^\xi(t)- d_{12}^\xi(t)\bigl(d_{22}^\xi(t)\bigr)\pinv d_{12}^\xi(t)^\transpose.
\]
\end{proposition}

\begin{remark}[Why the recursion is fast]
Each step consists of:
(i) a prediction of a mean and covariance using matrix multiplications,
(ii) one pseudo-inverse of the (conditional) observation covariance \(d_{22}^\xi(t)\),
(iii) a linear correction by the innovation \(\nu_{t+1}\).
\end{remark}

% ==========================================================
\section{Recovering the classical Kalman filter (M\"uller lectures notes)}

The lecture note (M\"uller) treats the standard linear-Gaussian state space model
\[
X_k = A X_{k-1} + B u_{k-1} + W_k,\qquad
Z_k = H X_k + R_k,
\]
with independent Gaussian noises \(W_k\sim \mathcal{N}(0,\Sigma_W)\), \(R_k\sim\mathcal{N}(0,\Sigma_R)\),
and information \(\F_k=\sigma(Z_1,\dots,Z_k)\).
In that setting \(d_{22}^\xi(t)\) is typically invertible and the pseudo-inverse becomes an inverse.

To match notations, set \(\theta_k=X_k\), \(\xi_k=Z_k\),
\(a_0(k-1)=B u_{k-1}\), \(a_1(k-1)=A\), \(A_0(k-1)=0\), \(A_1(k-1)=H\),
and choose matrices \(b_i,B_i\) so that
\(
b_1b_1^\transpose+b_2b_2^\transpose=\Sigma_W
\)
and
\(
B_1B_1^\transpose+B_2B_2^\transpose=\Sigma_R
\),
with \(b_1B_1^\transpose+b_2B_2^\transpose=0\) (uncorrelated state/measurement noise).
Then \eqref{eq:update_m}--\eqref{eq:update_Y} become the familiar Kalman filter:
\[
\hat X_k^- = A\hat X_{k-1}+B u_{k-1},\qquad
\Sigma_k^- = A\Sigma_{k-1}A^\transpose+\Sigma_W,
\]
and
\[
K_k = \Sigma_k^- H^\transpose\,(H\Sigma_k^-H^\transpose+\Sigma_R)^{-1},
\qquad
\hat X_k = \hat X_k^- + K_k(Z_k-H\hat X_k^-),
\qquad
\Sigma_k = \Sigma_k^- - K_k(H\Sigma_k^-H^\transpose+\Sigma_R)K_k^\transpose.
\]

% ==========================================================
\section{Example: 1D vehicle tracking (position \& velocity)}

In the lecture note example, the state is
\(
X_k=(p_k,v_k)^\transpose
\)
(position and velocity), with sampling interval \(\Delta>0\).
The (deterministic) kinematics under commanded acceleration \(u_k\) are
\[
p_k = p_{k-1}+\Delta v_{k-1}+\tfrac12 \Delta^2 u_{k-1},\qquad
v_k = v_{k-1}+\Delta u_{k-1}.
\]
Random acceleration noise is modeled by replacing \(u_{k-1}\) with \(u_{k-1}+E_{k-1}\),
where \(E_k\sim \mathcal{N}(0,\sigma_E^2)\) i.i.d.
This yields the state equation
\[
X_k =
\begin{pmatrix}
1 & \Delta\\
0 & 1
\end{pmatrix}
X_{k-1}+
\begin{pmatrix}
\tfrac12\Delta^2\\
\Delta
\end{pmatrix}
u_{k-1}
+W_k,
\quad
W_k=
\begin{pmatrix}
\tfrac12\Delta^2 E_{k-1}\\
\Delta E_{k-1}
\end{pmatrix}
\sim \mathcal{N}(0,\Sigma_W),
\]
with
\[
\Sigma_W=\sigma_E^2
\begin{pmatrix}
\tfrac14\Delta^4 & \tfrac12\Delta^3\\
\tfrac12\Delta^3 & \Delta^2
\end{pmatrix}.
\]
Only the position is observed with noise:
\[
Z_k = H X_k + R_k,\qquad H=\begin{pmatrix}1&0\end{pmatrix},
\qquad R_k\sim \mathcal{N}(0,\sigma_R^2).
\]
The filter then follows the classical recursion in the previous section.

% ==========================================================
\section{What to take away conceptually}

\begin{itemize}
\item The \textbf{core mathematical tool} is Gaussian conditioning (normal correlation). The update is
      always ``prior mean + gain $\times$ innovation'' and ``prior covariance minus explained part''.
\item The \textbf{generalization} in your course notes is to allow random/adaptive coefficients that are
      measurable w.r.t.\ the past observations, while preserving conditional Gaussianity.
\item The \textbf{pseudo-inverse} appears naturally when the conditional observation covariance is only
      positive semidefinite; the regression formulas remain valid without requiring invertibility.
\end{itemize}


